\chapter{Conclusion and Outlook}

Overall, the comparison of the wind estimation performance demonstrates that both the EKF and UKF provide comparable wind estimation performance for the investigated airship operating conditions. Although the UKF avoids the explicit model linearization required by the EKF, this does not result in a significant reduction in estimation error within the considered test range. The observed differences are comparatively small and depend on the individual scenario rather than consistently favoring one filter. Consequently, the EKF provides a suitable alternative to the UKF for the subsequent analyses, while offering the advantage of a lower computational effort associated with the propagation of a single linearized state model instead of multiple sigma points.





Outlook

The model-based approach proves to be a powerful tool for wind estimation given a good model of the system is available.  
The Dryden wind model does not only provide a mathematical description for the wind speeds acting at the center of gravity but also yields a description for the effect of the wind on the angular rates, which are caused by differences in wind magnitudes within the prevailing wind field.

Including the angular rates induced by the differences in wind speeds within the wind field can improve the performance as the airship is especially vulnerable to wind fields considering the airship's large side surfaces. 
Angular rates of the wind are described by transition functions of order 2 (for the u-direction) and order 3 (for v, w direction) resulting in a increase of one order for every direction compared to wind speeds.
It remains to investigate whether the gain in performance can justify the increased complexity.

The Wind estimation proved to be particulary reliable. This information can be used for advanced control and trajectory planning algorithms like MPC to increase the performance compared to the current PID setup.

Increase performance for hovering

Split up estimation in turbulence and mean wind to correctly use the mean wind CS for the Dryden turbulence.
